% ============================================================
% PRM393 - Section 2: Architecture & Foundation (Bilingual)
% ============================================================
\section{\bi{Kiến trúc ứng dụng di động}{Mobile System Architecture}}
\label{sec:arch}

\bi{
OmniCast Mobile tuân thủ triệt để nguyên lý phân tách trách nhiệm (Separation
of Concerns) và kiến trúc Clean Architecture được điều chỉnh cho hệ sinh thái
Flutter. Cấu trúc mã nguồn được phân chia thành bốn tầng độc lập:
Presentation, Logic (BLoC), Data và Core, bảo đảm tính độc lập khi kiểm thử
(testability), khả năng tái sử dụng mã nguồn và tính dễ bảo trì.
}{
OmniCast Mobile strictly adheres to the Separation of Concerns and Clean
Architecture principles adapted for the Flutter framework. The codebase is
organized into four distinct layers: Presentation, Logic (BLoC), Data, and
Core, ensuring high testability, modular reusability, and long-term
maintainability.
}

\begin{figure}[t]
\centering
\fitcol{\figMobileArchitecture}
\caption{\bi{Sơ đồ kiến trúc phân tầng 4 lớp của ứng dụng di động OmniCast:
tầng Presentation điều hướng qua GoRouter, tầng Logic với 12 BLoC/Cubit, tầng
Data giao tiếp đồng thời với REST API và SQLite ngoại tuyến.}{The 4-tier
layered architecture of OmniCast Mobile: Presentation routed by GoRouter,
Logic driven by 12 BLoC/Cubits, and Data communicating concurrently with
NestJS REST APIs and local SQLite storage.}}
\label{fig:mob_arch}
\end{figure}

\subsection{\bi{Phân tầng trách nhiệm}{Layered Organization}}

\bi{
Bốn tầng kiến trúc được tổ chức như sau:
}{
The four architectural tiers are defined as follows:
}

\begin{enumerate}[leftmargin=*,itemsep=2pt]
  \item \textbf{Presentation Layer:} \bi{Gồm 31 màn hình (screens) và 22
  thành phần giao diện tái sử dụng (widgets). Tầng này hoàn toàn không chứa
  logic nghiệp vụ phức tạp mà chỉ chịu trách nhiệm hiển thị trạng thái và
  phát sinh các sự kiện người dùng (User Events).}{%
  Encompasses 31 screens and 22 reusable widgets. This tier is completely free
  of business computations, focusing solely on state rendering and event
  dispatching.}

  \item \textbf{Logic Layer (BLoC / Cubit):} \bi{Gồm 12 module điều phối
  nghiệp vụ. BLoC tiếp nhận sự kiện từ Presentation, giao tiếp với tầng Data
  qua các Repository và phát ra các trạng thái mới (States) thông qua các
  Stream bất đồng bộ.}{%
  Hosts 12 state-management modules. BLoCs accept incoming events from UI
  widgets, invoke domain repositories, and emit reactive state streams.}

  \item \textbf{Data Layer:} \bi{Chứa các Repository, Data Sources và Model
  dữ liệu. Tầng này trừu tượng hóa nguồn gốc dữ liệu: tự động quyết định đọc
  từ bộ nhớ đệm SQLite hay gọi mạng từ xa qua Dio Client.}{%
  Encapsulates Repositories, Data Sources, and Data Transfer Objects (DTOs).
  It abstracts data origins, transparently querying SQLite or fetching
  remotely via Dio.}

  \item \textbf{Core Layer:} \bi{Cung cấp các tiện ích dùng chung: bộ định
  tuyến \texttt{GoRouter}, bộ tiêm phụ thuộc \texttt{get\_it}, cấu hình
  giao diện (theme), chuỗi bản địa hóa (i18n), dịch vụ thông báo cục bộ và
  quản lý token bảo mật.}{%
  Supplies cross-cutting concerns: \texttt{GoRouter} navigation definitions,
  \texttt{get\_it} dependency injection, UI theme styling, i18n
  localization, and secure token storage.}
\end{enumerate}

\subsection{\bi{Bộ tiêm phụ thuộc (Dependency Injection)}{Dependency Injection with GetIt}}

\bi{
Toàn bộ các thực thể Repository, Data Source, HTTP Client và BLoC được khởi
tạo và quản lý vòng đời thông qua thư viện \texttt{get\_it}. Trong tệp
\texttt{injection\_container.dart}, cơ chế đăng ký được phân định rõ ràng
giữa Lazy Singleton (cho các dịch vụ mạng, cơ sở dữ liệu và repository) và
Factory (cho các BLoC cần tạo mới theo vòng đời màn hình):
}{
All Repositories, Data Sources, Network Clients, and BLoC instances have
their lifecycles managed via the \texttt{get\_it} service locator. In
\texttt{injection\_container.dart}, registrations are strictly partitioned
between Lazy Singletons (for stateless services, DB helpers, and repositories)
and Factories (for transient BLoC instances tied to screen lifecycles):
}

\begin{lstlisting}[language=Dart, caption={\bi{Cấu hình tiêm phụ thuộc trong \texttt{injection\_container.dart}.}{Dependency injection configuration in \texttt{injection\_container.dart}.}}]
final sl = GetIt.instance;

Future<void> init() async {
  // BLoCs (Factory)
  sl.registerFactory(() => AuthBloc(sl()));
  sl.registerFactory(() => WatchlistBloc(sl()));
  sl.registerFactory(() => EpgBloc(sl()));

  // Repositories (Lazy Singleton)
  sl.registerLazySingleton<WatchlistRepository>(
    () => WatchlistRepositoryImpl(
      remoteDataSource: sl(),
      localDataSource: sl(),
    ),
  );

  // Core & Drivers
  sl.registerLazySingleton(() => DioClient(sl()));
  sl.registerLazySingleton(() => DatabaseHelper());
}
\end{lstlisting}

\subsection{\bi{Tầng mạng và xoay vòng Token bảo mật}{Networking and JWT Token Rotation}}

\bi{
Giao tiếp HTTP được thực hiện qua thư viện \texttt{Dio} với lớp đánh chặn
\texttt{AuthInterceptor}. Khi ứng dụng nhận mã lỗi \texttt{HTTP 401
Unauthorized}, bộ đánh chặn sẽ tự động kích hoạt tiến trình làm mới token
thông qua Refresh Token lưu trong \texttt{flutter\_secure\_storage}:
}{
Remote HTTP communications are powered by \texttt{Dio} enriched with a custom
\texttt{AuthInterceptor}. Upon intercepting an \texttt{HTTP 401 Unauthorized}
response, the interceptor transparently initiates a background token refresh
using the refresh token persisted in \texttt{flutter\_secure\_storage}:
}

\begin{enumerate}[leftmargin=*,itemsep=2pt]
  \item \bi{Tạm dừng hàng đợi các yêu cầu đang chờ (request queue lock).}{%
  Locks the active request queue to prevent cascading 401 errors.}
  \item \bi{Gửi yêu cầu làm mới tới \texttt{/api/v1/auth/refresh} với Refresh
  Token hợp lệ.}{%
  Dispatches a refresh request to \texttt{/api/v1/auth/refresh} with the
  existing refresh token.}
  \item \bi{Lưu trữ cặp Access/Refresh token mới và phát lại các yêu cầu bị
  tạm dừng với Bearer token mới.}{%
  Persists the refreshed token pair and replays pending requests with the
  new Bearer token.}
  \item \bi{Nếu Refresh Token hết hạn hoặc bị thu hồi, chuyển hướng người
  dùng về màn hình Đăng nhập một cách an toàn.}{%
  If the refresh token itself is invalid or revoked, safely clears the session
  and redirects the user to the Login screen.}
\end{enumerate}
